% Part: turing-machines
% Chapter: undecidability
% Section: unsolvability-decision-problem

\documentclass[../../../include/open-logic-section]{subfiles}

\begin{document}

\olfileid{tur}{und}{uns}
\olsection{નિર્ણય સમસ્યા ઉકેલી શકાતી નથી}

\begin{thm}
\ollabel{thm:decision-prob}
નિર્ણય સમસ્યા ઉકેલી શકાતી નથી: એવું કોઈ ટ્યુરિંગ
મશીન~$D$ નથી કે જેની ટેપ પર પ્રથમ-ક્રમના તર્કનું
!!a{sentence}~$!B$ ઇનપુટ તરીકે હોય ત્યારે~$D$ છેવટે
અટકે અને~$!B$ પ્રામાણ્ય ધરાવે ત્યારે અને માત્ર ત્યારે
જ~$1$, અન્યથા~$0$ આઉટપુટ આપે.
\end{thm}

\begin{proof}
માનો કે નિર્ણય સમસ્યા ઉકેલી શકાતી હોય, એટલે એવું
ટ્યુરિંગ મશીન~$D$ હોય. તો અટકવાની સમસ્યા આ પ્રમાણે
ઉકેલી શકીએ. ટ્યુરિંગ મશીન~$M_e$નો સૂચક~$e$ અને
ઇનપુટ~$w$ આપતાં, તેને અનુરૂપ !!{sentence}
$!T(M_e, w) \lif !E(M_e, w)$ ગણતું અને ટેપના સૌથી
ડાબા ખાનાને વાંચતું અટકતું ટ્યુરિંગ મશીન~$E$ બનાવીએ.
પછી $E \concat D$ મશીન ઇનપુટ~$e$ અને~$w$ માટે
પહેલાં $!T(M_e, w) \lif !E(M_e, w)$ ગણે અને ત્યાર પછી
તે ઇનપુટ પર નિર્ણય સમસ્યાનું મશીન~$D$ ચલાવે.
$!T(M_e, w) \lif !E(M_e, w)$ પ્રામાણ્ય ધરાવે ત્યારે
અને માત્ર ત્યારે $D$ આઉટપુટ~$1$ સાથે અટકે છે;
અન્યથા આઉટપુટ~$0$ આપે છે.
\olref[ver]{lem:halt-if-valid} અને
\olref[ver]{lem:valid-if-halt} મુજબ
$!T(M_e, w) \lif !E(M_e, w)$ પ્રામાણ્ય
ધરાવે ત્યારે અને માત્ર ત્યારે~$M_e$ ઇનપુટ~$w$ પર
અટકે છે. એટલે ઇનપુટ~$e$ અને~$w$ પર $E\concat D$
એવું મશીન બને કે~$M_e$ ઇનપુટ~$w$ પર અટકે ત્યારે
અને માત્ર ત્યારે આઉટપુટ~$1$ સાથે, અન્યથા આઉટપુટ~$0$
સાથે અટકે. બીજા શબ્દોમાં, $E \concat D$ અટકવાની
સમસ્યા ઉકેલી નાખે. પરંતુ
\olref[hal]{thm:halting-problem}થી આપણે જાણીએ છીએ
કે આવું કોઈ ટ્યુરિંગ મશીન હોઈ શકતું નથી.
\end{proof}

\begin{cor}\ollabel{cor:undecidable-sat}%
પ્રથમ-ક્રમના તર્કનું મનસ્વી !!{sentence} સંતોષ્ય છે કે
નહીં તે નક્કી કરી શકાતું નથી.
\end{cor}

\begin{proof}
  માનો કે સંતોષ્યતા ટ્યુરિંગ મશીન~$S$ વડે નક્કી કરી
  શકાતી હોય. તો નિર્ણય સમસ્યા આ રીતે ઉકેલી શકીએ:
  !!a{sentence}~$!B$ ઇનપુટ તરીકે આપેલું હોય, તો
  ટેપ પરનું~$!B$ એક ખાનું જમણી તરફ ખસેડો.
  \sourcecorrection{OLUN-005}{સ્થિર અંગ્રેજી મૂળ આ
  વાક્યમાં ઇનપુટને $B$ કહે છે; આસપાસની સાબિતીમાં
  તે જ વાક્ય $!B$ છે, તેથી અહીં ચિહ્ન એકરૂપ કર્યું છે.}
  ખાનું~$1$ પર પાછા જઈ~$\lnot$ ચિહ્ન લખો.

  હવે ટ્યુરિંગ મશીન~$S$ ચલાવો. ટેપ પર આઉટપુટ~$1$
  (જો $\lnot !B$ સંતોષ્ય હોય) અથવા~$0$ (જો
  $\lnot !B$ અસંતોષ્ય હોય) સાથે તે છેવટે અટકે છે.
  ખાનું~$1$ પર~$\TMstroke$ હોય, તો તેને ભૂંસી નાખો;
  ખાનું~$1$ ખાલી હોય, તો ત્યાં~$\TMstroke$ લખો,
  અને પછી અટકો.

  આ ટ્યુરિંગ મશીન હંમેશાં અટકે છે અને~$\lnot !B$
  અસંતોષ્ય હોય ત્યારે અને માત્ર ત્યારે આઉટપુટ~$1$,
  અન્યથા~$0$ આપે છે. $!B$ પ્રામાણ્ય ધરાવે ત્યારે
  અને માત્ર ત્યારે~$\lnot !B$ અસંતોષ્ય છે. એટલે
  મશીન~$!B$ પ્રામાણ્ય ધરાવે ત્યારે અને માત્ર ત્યારે
  આઉટપુટ~$1$, અન્યથા~$0$ આપે છે; આમ તે નિર્ણય
  સમસ્યા ઉકેલી નાખે.
\end{proof}

\begin{explain}
એટલે ``$!B$ પ્રથમ-ક્રમના તર્કનું પ્રમાણભૂત !!{sentence}
છે?'' એ પ્રશ્નનો હંમેશાં સાચો ``હા'' અથવા ``ના''
જવાબ આપતું ટ્યુરિંગ મશીન નથી. પરંતુ સાચો જવાબ
``હા'' હોય ત્યારે તે આપતું, અને જવાબ ``ના'' હોય તો
અટકતું જ ન હોય એવું ટ્યુરિંગ મશીન \emph{છે}.
આ વાત પ્રથમ-ક્રમના તર્કની યથાર્થતા અને પૂર્ણતાના
પ્રમેયોમાંથી તથા !!{derivation}sની અસરકારક પરિગણના
કરી શકાય છે એ હકીકતમાંથી ફલિત થાય છે.
\end{explain}

\begin{thm}
  \ollabel{thm:valid-ce}%
  પ્રથમ-ક્રમનાં !!{sentence}sનું પ્રામાણ્ય અર્ધનિર્ણેય
  છે: એવું ટ્યુરિંગ મશીન~$E$ છે કે જેની ટેપ પર
  પ્રથમ-ક્રમનું !!a{sentence}~$!B$ ઇનપુટ તરીકે હોય,
  તો~$!B$ પ્રામાણ્ય ધરાવે ત્યારે અને માત્ર ત્યારે
  $E$ છેવટે આઉટપુટ~$1$ સાથે અટકે છે; અન્યથા તે
  અટકતું નથી.
\end{thm}

\begin{proof}
  પ્રથમ-ક્રમના તર્કની બધી શક્ય !!{derivation}s એક
  પછી એક અસરકારક વિધિ વડે ઉત્પન્ન કરી શકાય છે.
  મશીન~$E$ એમ કરે છે અને~$\Proves !B$ બતાવતી
  !!a{derivation} મળે ત્યારે આઉટપુટ~$1$ સાથે અટકે છે.
  યથાર્થતાના પ્રમેયથી, $E$ આઉટપુટ~$1$ સાથે અટકે,
  તો~$\Entails !B$ છે. પૂર્ણતાના પ્રમેયથી,
  $\Entails !B$ હોય તો~$\Proves !B$ બતાવતી
  !!a{derivation} છે. $E$ બધી શક્ય !!{derivation}s
  પદ્ધતિસર ઉત્પન્ન કરતું હોવાથી તે
  $\Proves !B$ બતાવતી વ્યૂત્પત્તિ છેવટે
  શોધશે અને આઉટપુટ~$1$
  સાથે અટકશે.
\end{proof}

\end{document}
